% ===================================================================
%  TABEL Nr. 6 — Het huis van binnen, het meubilair, de voeding, de
%  verlichting.
%  Traduction 2026 d'apres texto/io, controlee sur texto/fr et
%  texto/es. Consigne : voir texto/nl/10-tabel-01.tex.
%
%  Deux notes, toutes deux marquees « (*) », et deux appels. L'ordre
%  les departage.
%
%  LA FEMME DE CHAMBRE PREND SON (16), contre le fac-simile ido qui
%  ecrit (6) — son propre numero du savon. Ecart declare dans
%  outils/renvoji.py.
% ===================================================================

%%K t06-apar-1 apar
\VUsaut{6.0mm}
\VUtitre{15.0pt}{60}{TABEL Nr. 6}
\VUsaut{1.4mm}
\VUtitre{11.4pt}{0}{Het huis van binnen, het meubilair, de}
\VUsaut{0.4mm}
\VUtitre{11.4pt}{0}{voeding, de verlichting.}
\VUsaut{2.2mm}

%%K t06-apar-2 apar
Het huis is afgebouwd. De oom van Jan is in zijn nieuwe woning geïnstalleerd, waarvan wij de
indeling kennen. Laten wij nu zien hoe hij zijn huis gemeubileerd heeft. Wij zullen eerst
verkennen:

%%K t06-tit-1 sub
\VUpk{10.2pt}{40}{De slaapkamer.}
\VUsaut{3.0mm}

%%K t06-c1-01-1 p
1. --- Wij komen ongelegen, want de kamenier is bezig de \VUgras{kamer} schoon te maken.

%%K t06-c1-02-1 p
2. --- Op de \VUgras{wastafel} \textsuperscript{(1)} liggen hier en daar verschillende
voorwerpen waarmee men zich wast, zich kamt, kortom zijn toilet maakt. Het zijn de
\VUgras{waskom} \textsuperscript{(2)} en de
\VUgras{waterkan} \textsuperscript{(3)}, de \VUgras{haarborstel}
\textsuperscript{(4)}, de \VUgras{kam} \textsuperscript{(5)}, de \VUgras{zeep}
\textsuperscript{(6)} en de \VUgras{spons} \textsuperscript{(7)}, ten slotte een
\VUgras{flesje} \textsuperscript{(8)} voor het vloeibare tandmiddel.

%%K t06-c1-03-1 p
3. --- Op de grond, voor de wastafel staat de \VUgras{emmer} \textsuperscript{(9)} om het
vuile water op te vangen.

%%K t06-c1-04-1 p
4. --- Wij zien in de \VUgras{voorkamer} \textsuperscript{(10)} enkele
\VUgras{kapstokhaken} \textsuperscript{(13)} en twee \VUgras{paraplu's}
\textsuperscript{(11)} in de \VUgras{parapluhouder} \textsuperscript{(12)}.

%%K t06-c1-05-1 p
5. --- Aan de andere kant van de deur staat de \VUgras{spiegelkast} \textsuperscript{(14)} die
het \VUgras{linnengoed} \textsuperscript{(15)} bevat.

%%K t06-c1-06-1 p
6. --- Laten wij ervan profiteren dat de \VUgras{kamenier} \textsuperscript{(16)} het
\VUgras{bed} \textsuperscript{(17)} opmaakt, om de verschillende delen ervan te verkennen. Het
\VUgras{ledikant} \textsuperscript{(20)} bevat een goede \VUgras{springveren matras}
\textsuperscript{(21)} waarop Justine dadelijk een wollen \VUgras{matras}
\textsuperscript{(22)} zal leggen. Daarna zal zij van de stoelen de twee
\VUgras{lakens} \textsuperscript{(23)} en de \VUgras{deken}
\textsuperscript{(24)} nemen. Dan zal zij op het hoofdeinde het
\VUgras{dwarskussen} \textsuperscript{(25)} en het \VUgras{hoofdkussen}
\textsuperscript{(26)} leggen die met het \VUgras{dekbed} \textsuperscript{(27)} op het
\VUgras{beddenkleed} \textsuperscript{(32)} liggen. Zij gebruikt een plumeau om de
\VUgras{gordijnen} \textsuperscript{(19)} en de \VUgras{bedhemel} \textsuperscript{(18)} af te
stoffen, en zo is een deel van het werk klaar.

%%K t06-c1-07-1 p
7. --- Dan zal zij het \VUgras{nachtkastje} \textsuperscript{(28)} wat moeten opruimen, de
\VUgras{wekker} \textsuperscript{(29)} opnieuw opwinden, het
\VUgras{nachtlampje} \textsuperscript{(30)} klaarmaken, de \VUgras{kaars}
\textsuperscript{(31)} wegnemen om de kandelaar schoon te maken.

%%K t06-tit-2 sub
\VUsaut{5.0mm}
\VUpk{10.2pt}{40}{Het werkmandje en het haardgerei.}
\VUsaut{3.0mm}

%%K t06-c1-08-1 p
8. --- Het \VUgras{werkmandje} \textsuperscript{(34)} naast het \VUgras{krukje}
\textsuperscript{(33)} moet ook heel nodig opgeruimd worden. Zij zal de
\VUgras{borduurwerken} \textsuperscript{(35)} moeten opvouwen, in het
\VUgras{werktafeltje} \textsuperscript{(38)} de \VUgras{vingerhoed}, de
\VUgras{draad} \textsuperscript{(36)}, de \VUgras{naalden}
\textsuperscript{(37)} en de \VUgras{schaar} \textsuperscript{(39)} opbergen en met de
\VUgras{sleutel} \textsuperscript{(40)} het deksel van het tafeltje sluiten.

%%K t06-c1-09-1 p
9. --- Op het \VUgras{eenpotige tafeltje} \textsuperscript{(41)} in het midden zal Justine het
water van de \VUgras{karaf} \textsuperscript{(42)} verversen. Zij zal \VUgras{suiker} in het
\VUgras{suikerpotje} \textsuperscript{(43)} doen, het \VUgras{suikertangetje}
\textsuperscript{(44)} schoonmaken en de
\VUgras{kleerborstel} \textsuperscript{(46)} wegnemen.

%%K t06-c1-10-1 p
10. --- De rechterkant van de kamer zal minder werk van haar vragen, want de
\VUgras{chaise longue} \textsuperscript{(47)} en haar \VUgras{kussen}
\textsuperscript{(48)} staan op hun juiste plaats, en, omdat het geen koud weer is, is niets
verplaatst tussen het gerei van de \VUgras{schoorsteen} \textsuperscript{(49)}. De
\VUgras{schop} \textsuperscript{(50)}, de
\VUgras{tang} \textsuperscript{(51)}, de \VUgras{haardijzers}
\textsuperscript{(55)}, het \VUgras{ashek} \textsuperscript{(56)} zijn schoon; het
\VUgras{vuurscherm} \textsuperscript{(54)} is niet verplaatst en de
\VUgras{houtbak} \textsuperscript{(52)} is goed voorzien van
\VUgras{brandhout} \textsuperscript{(53)}.

%%K t06-noto-1 noto
\VUnotes{143.2mm}{%
(*) Babo = kleuter, het E. baby, F. bébé, I. bambino.%
}

%%K t06-noto-2 noto
\VUnotes{143.2mm}{%
(*) Lambrequino = F. lambrequin, S. lambrequines, Port. lambrequins, zelfs D. E. lambrequins,
dus D. E. F. P. S.%
}

%%K t06-c1-11-1 p
11. --- Toch zal zij de \VUgras{pendule} \textsuperscript{(57)}, de
\VUgras{kandelabers} \textsuperscript{(58)} en de \VUgras{prullenbakjes}
\textsuperscript{(59)} moeten afstoffen, ten slotte de \VUgras{spiegel} \textsuperscript{(60)}
afvegen.

%%K t06-c1-12-1 p
12. --- Wat de \VUgras{wieg} \textsuperscript{(61)} van het kindje (van de babo (*)) betreft,
die staat al klaar.

%%K t06-tit-3 sub
\VUsaut{5.0mm}
\VUpk{10.2pt}{40}{De salon.}
\VUsaut{3.0mm}

%%K t06-c2-01-1 p
1. --- Nu gaan we de salon binnen. Het is een heel mooie \VUgras{kamer}. De schoorsteen, die
in de rechterhoek staat is versierd met een fijn
\VUgras{beeldje} \textsuperscript{(1)}, met twee mooie \VUgras{vazen}
\textsuperscript{(3)} en gedrapeerd met een fluwelen \VUgras{schoorsteenval}
\textsuperscript{(2)} (*).

%%K t06-c2-02-1 p
2. --- Om te beletten dat de vonken van de haard het tapijt in brand steken, heeft men voor de
haard een koperen \VUgras{vonkenscherm} \textsuperscript{(4)} gezet.

%%K t06-c2-03-1 p
3. --- Dichtbij staat een sierlijk \VUgras{schrijftafeltje} \textsuperscript{(5)} voor dames
open. Daaraan schrijft de \VUgras{vrouw des huizes} \textsuperscript{(23)} haar brieven.

%%K t06-c2-04-1 p
4. --- Het raam is voorzien van \VUgras{vitrages} \textsuperscript{(6)} met
\VUgras{koorden} \textsuperscript{(8)} die aan \VUgras{haken}
\textsuperscript{(7)} bevestigd zijn en van stoffen overgordijnen
\VUgras{die bovenaan gedrapeerd} zijn \textsuperscript{(9)}.

%%K t06-c2-05-1 p
5. --- In de vensternis bevat een \VUgras{sierpot} \textsuperscript{(11)} een groene
\VUgras{plant} \textsuperscript{(10)}, een prachtige palmboom waarvan de bladeren zich bijna
tot aan de \VUgras{petroleumlamp} \textsuperscript{(12)} uitstrekken.

%%K t06-c2-06-1 p
6. --- De \VUgras{lampenkap} \textsuperscript{(13)} van de lamp is geschikt door Maria, het
bekoorlijke \VUgras{meisje} \textsuperscript{(14)} dat aan de piano \textsuperscript{(15)}
zit. Heel recht op het \VUgras{krukje} \textsuperscript{(16)} zittend, studeert zij een
muziekstuk. Zij is heel bekwaam en zorgvuldig, zoals haar \VUgras{muziekkast}
\textsuperscript{(17)} bewijst die volmaakt opgeruimd is.

%%K t06-tit-4 sub
\VUsaut{5.0mm}
\VUpk{10.2pt}{40}{De deugniet.}
\VUsaut{3.0mm}

%%K t06-c2-07-1 p
7. --- Haar broer, Paul, zoekt een \VUgras{boek} \textsuperscript{(19)} in de
\VUgras{boekenkast} \textsuperscript{(18)}. Hij is een \VUgras{jongen}
\textsuperscript{(20)}, woelig en onuitstaanbaar. Terwijl hij zich aan de boekenkast optrekt
om het boek te bereiken dat te hoog staat, loopt hij gevaar de \VUgras{buste}
\textsuperscript{(21)} te laten vallen die erboven staat.

%%K t06-c2-08-1 p
8. --- Hij maakt van de gelegenheid gebruik om die dwaasheid te doen, omdat zijn moeder bezig
is te praten met een vriendin, een \VUgras{dame} \textsuperscript{(22)} die enkele dagen bij
haar kwam doorbrengen. De moeder zit op een \VUgras{canapé} \textsuperscript{(24)} naast de
\VUgras{leunstoel} \textsuperscript{(25)}, en haar vriendin zit op de \VUgras{stoel}
\textsuperscript{(27)} waarvan men alleen de \VUgras{rugleuning} \textsuperscript{(26)} ziet.

%%K t06-c2-09-1 p
9. --- De grote \VUgras{lijst} \textsuperscript{(30)} die aan de muur hangt bevat het
\VUgras{portret} \textsuperscript{(29)} van een bloedverwante. In het
\VUgras{album} \textsuperscript{(32)} dat op de tafel ligt zijn de
\VUgras{foto's} \textsuperscript{(33)} van de overige familieleden verzameld. De
kinderen hebben het zojuist doorgebladerd, want de \VUgras{stoelen} \textsuperscript{(31)}
staan nog om de tafel gegroepeerd.

%%K t06-c2-10-1 p
10. --- De porseleinen \VUgras{Chinese vaas} \textsuperscript{(34)}, die naast de
\VUgras{briefopener} \textsuperscript{(35)} staat werd aan Maria voor haar
naamdag geschonken, en het \VUgras{kamerscherm} \textsuperscript{(36)} dat de hoek van de
kamer siert werd door haar oom uit China meegebracht.

%%K t06-c2-11-1 p
11. --- Opgevrolijkt door de prachtige \VUgras{schilderijen} \textsuperscript{(37)} die aan de
\VUgras{wand} \textsuperscript{(42)} hangen, verlicht door de \VUgras{plafondluchter}
\textsuperscript{(38)}, waarvan de
\VUgras{elektrische lampen} \textsuperscript{(39)} 's avonds vrolijk schijnen,
ziet die salon er heel aangenaam uit. Het zachte \VUgras{tapijt} \textsuperscript{(40)} dat op
het parket ligt, en de zo gemakkelijk te gebruiken
\VUgras{elektrische bel} \textsuperscript{(41)}, voltooien het comfort ervan.

%%K t06-tit-5 sub
\VUsaut{3.6mm}
\VUpk{10.2pt}{40}{De eetkamer.}
\VUsaut{3.0mm}

%%K t06-c3-01-1 p
1. --- Het dineruur heeft geslagen. Het hele gezin, behalve Maria, is om de tafel verzameld.
De eetkamer wordt aangenaam verwarmd door de uitstekende
\VUgras{kachel} \textsuperscript{(1)} van aardewerk, met een dunne
\VUgras{pijp} \textsuperscript{(2)}.

%%K t06-c3-02-1 p
2. --- Die kamer bevat niet veel meubels. Langs de muur hangt, boven het
\VUgras{bankje} \textsuperscript{(4)}, een sierlijk \VUgras{fonteintje}
\textsuperscript{(3)}. Vlakbij staat het \VUgras{buffet} \textsuperscript{(5)} waarop men de
koude spijzen, de nagerechtschotels, en het verschillende tafelgerei zet dat men niet dadelijk
nodig heeft.

%%K t06-c3-03-1 p
3. --- Zo zien wij op de bovenste plank een \VUgras{gebraden kip} \textsuperscript{(7)}, een
\VUgras{vis} \textsuperscript{(8)} en een
\VUgras{schaal} \textsuperscript{(10)} met \VUgras{vruchten}
\textsuperscript{(9)}. Daaronder liggen de \VUgras{kaas} \textsuperscript{(11)}, het
\VUgras{brood} \textsuperscript{(12)}, het
\VUgras{olie-en-azijnstel} \textsuperscript{(13)} dat alles bevat wat men nodig
heeft om de sla aan te maken: de olie, de azijn, het \VUgras{zout} \textsuperscript{(14)} en
de \VUgras{peper} \textsuperscript{(15)}. Ten slotte staat op de onderste plank een
\VUgras{taart} \textsuperscript{(17)} naast het
\VUgras{mosterdpotje} \textsuperscript{(16)}.

%%K t06-c3-04-1 p
4. --- De klok van de eetkamer, die men \VUgras{wandklok} \textsuperscript{(20)} noemt, heeft
een heel bijzondere vorm. Zij is in een houten lijst gevat die bovendien een
\VUgras{barometer} \textsuperscript{(19)} en een \VUgras{thermometer} \textsuperscript{(18)}
bevat.

%%K t06-tit-6 sub
\VUsaut{4.6mm}
\VUpk{10.2pt}{40}{Het opdienen van het diner.}
\VUsaut{3.0mm}

%%K t06-c3-05-1 p
5. --- Tegen alle muren van de kamer zijn \VUgras{houten panelen} \textsuperscript{(21)} van
gewaste eik aangebracht. Een stoffen
\VUgras{deurgordijn} \textsuperscript{(22)} sluit voldoende de deur af die
toegang geeft tot de veranda. Daar zal het gezin na het diner de koffie gaan drinken. Al staan
de \VUgras{kopjes} \textsuperscript{(24)} en de
\VUgras{schoteltjes} \textsuperscript{(25)} klaar op het \VUgras{dientafeltje}
\textsuperscript{(23)}.

%%K t06-c3-06-1 p
6. --- Boven de tafel is aan het plafond een \VUgras{hanglamp} \textsuperscript{(26)}
bevestigd waarvan de heel heldere lamp 's avonds de hele eetkamer verlicht.

%%K t06-c3-07-1 p
7. --- Laten wij nu de \VUgras{eettafel} \textsuperscript{(27)} zelf verkennen. Toen het gezin
binnenkwam, stond het couvert klaar. Op het \VUgras{tafellaken} \textsuperscript{(28)} was op
de plaats van elke disgenoot een \VUgras{bord} \textsuperscript{(33)}, een \VUgras{servet}
\textsuperscript{(29)}, een
\VUgras{lepel} \textsuperscript{(34)}, een \VUgras{vork}
\textsuperscript{(35)}, een \VUgras{mes} \textsuperscript{(36)} en een
\VUgras{glas} \textsuperscript{(32)} gelegd. Er waren ook \VUgras{flessen}
\textsuperscript{(30)} voor de wijn en een \VUgras{karaf} \textsuperscript{(31)} voor het
water.

%%K t06-c3-08-1 p
8. --- De \VUgras{bediende} \textsuperscript{(39)} heeft eerst de
\VUgras{soepterrine} \textsuperscript{(37)} gebracht, waarin nog wat soep
dampt. Nu dient hij de eerste \VUgras{schotel} \textsuperscript{(38)} op, een vleesschotel.
Daarna zal hij die opdienen welke wij al genoemd hebben, ten slotte, na de
\VUgras{nagerechtschotel}, zal hij ervan profiteren dat zijn meesters de veranda zullen zijn
binnengegaan, om het tafelgerei weg te nemen en de eetkamer weer op te ruimen.

%%K t06-c3-08-2 p
\textit{Nota.} --- \textit{De gangbare woorden betreffende de voeding worden hier heel beknopt
aangegeven. De lijst zal aangevuld worden op de twee tabellen « Het Hotel » nr. 13) en « De
Markt » (nr. 15).}

%%K t06-tit-7 sub
\VUsaut{4.6mm}
\VUpk{10.2pt}{40}{De keuken.}
\VUsaut{3.0mm}

%%K t06-c4-01-1 p
1. --- In de keuken, die wij door de halfopen deur beginnen te bekijken, zien wij een
schilderachtige wanorde. Voor de \VUgras{haard} \textsuperscript{(1)} waarin een \VUgras{vuur}
\textsuperscript{(3)} knettert, staat een
\VUgras{braadspit} \textsuperscript{(5)}. Een mooi \VUgras{gebraad}
\textsuperscript{(7)} is \VUgras{aan het spit geregen} \textsuperscript{(6)} en gaart verder.

%%K t06-c4-02-1 p
2. --- De \VUgras{blaasbalg} \textsuperscript{(8)} en de \VUgras{kolenschop}
\textsuperscript{(9)}, die men zojuist gebruikt heeft, zijn op de grond blijven liggen evenals
het \VUgras{brandhout} \textsuperscript{(10)}.

%%K t06-c4-03-1 p
3. --- De schoorsteenmantel is opgeruimd; de \VUgras{lantaarn} \textsuperscript{(11)}, het
\VUgras{luciferdoosje} \textsuperscript{(13)}, waarin de \VUgras{lucifers}
\textsuperscript{(12)} zitten, de
\VUgras{kandelaar} \textsuperscript{(14)} en het \VUgras{blakertje}
\textsuperscript{(15)} met een \VUgras{kaars} \textsuperscript{(16)} staan op hun plaats. Maar
de grote \VUgras{kolenemmer} \textsuperscript{(18)}, vol
\VUgras{steenkool} \textsuperscript{(17)} die de \VUgras{keukenmeid}
\textsuperscript{(23)} zojuist in de kelder gevuld heeft, is midden in de keuken blijven
staan.

%%K t06-c4-04-1 p
4. --- Werkelijk, de arme vrouw moet zich haasten opdat het middagmaal op tijd klaar zou zijn.
Nu staat zij bij het \VUgras{keukenfornuis} \textsuperscript{(19)} en roert een \VUgras{saus}
\textsuperscript{(24)} die niet te veel mag aanbranden.

%%K t06-c4-05-1 p
5. --- In de \VUgras{oven} \textsuperscript{(20)} bakt een taart, die zij voor het vieruurtje
van de kinderen heeft klaargemaakt. Boven het keukenfornuis hangen de \VUgras{pollepel}
\textsuperscript{(21)} en de
\VUgras{schuimspaan} \textsuperscript{(22)}.

%%K t06-tit-8 sub
\VUsaut{4.0mm}
\VUpk{10.2pt}{40}{Het vaatwerk.}
\VUsaut{3.0mm}

%%K t06-c4-06-1 p
6. --- Het \VUgras{keukengerei} \textsuperscript{(25)}, dat achterin de kamer hangt, is heel
schoon. De mooie koperen \VUgras{kastrollen} \textsuperscript{(26)}, de \VUgras{melkkan}
\textsuperscript{(27)}, het
\VUgras{zeefje} \textsuperscript{(28)} en de \VUgras{rasp}
\textsuperscript{(29)} zijn zorgvuldig schoongemaakt.

%%K t06-c4-07-1 p
7. --- De \VUgras{zoutdoos} \textsuperscript{(30)} staat op de vaatkast naast de
\VUgras{theepot} \textsuperscript{(31)}, en de \VUgras{grote oliefles}
\textsuperscript{(32)}. De \VUgras{lade} \textsuperscript{(33)} bevat waarschijnlijk het
eetgerei en de keukenmessen.

%%K t06-c4-08-1 p
8. --- De \VUgras{bezem} \textsuperscript{(34)} en de \VUgras{plumeau} \textsuperscript{(35)}
staan in een hoek. De keukenmeid heeft zojuist het
\VUgras{hakblok} \textsuperscript{(36)} gebruikt, zij heeft het bij het raam
laten staan.

%%K t06-c4-09-1 p
9. --- Een \VUgras{handdoek} \textsuperscript{(37)} hangt aan de muur, naast de
\VUgras{trechter} \textsuperscript{(38)}. In dat deel van de keuken zijn de
\VUgras{rooster} \textsuperscript{(39)}, het \VUgras{hakmes}
\textsuperscript{(40)} en de \VUgras{hakplank} \textsuperscript{(41)} weggeborgen. Daarna
staan, boven het hakmes, op een \VUgras{plank} \textsuperscript{(42)}, \VUgras{potten}
\textsuperscript{(43)}, de
\VUgras{kookpot} \textsuperscript{(44)} met zijn \VUgras{deksel}
\textsuperscript{(45)} en de \VUgras{marmiet} \textsuperscript{(46)}. De
\VUgras{koekenpan} \textsuperscript{(47)} hangt dichtbij.

%%K t06-c4-10-1 p
10. --- Alleen de koperen \VUgras{kraan} \textsuperscript{(48)} werpt een glans in die hoek.
De \VUgras{gootsteen} \textsuperscript{(49)} waarop de dikke
\VUgras{kruik} \textsuperscript{(50)} staat, is waarschijnlijk heel handig met
zijn kleine kast, waarin men het \VUgras{azijnvaatje} \textsuperscript{(51)} bewaart dat van
zijn \VUgras{kraantje} \textsuperscript{(52)} voorzien is, de
\VUgras{afvalbak} \textsuperscript{(53)} en de \VUgras{vuile borden}
\textsuperscript{(54)}.

%%K t06-tit-9 sub
\VUsaut{4.0mm}
\VUpk{10.2pt}{40}{Andere voorwerpen in de keuken.}
\VUsaut{3.0mm}

%%K t06-c4-11-1 p
11. --- De \VUgras{tobbe} \textsuperscript{(55)} waarin men de
\VUgras{groenten} \textsuperscript{(58)} wast en de gietijzeren \VUgras{ketel}
\textsuperscript{(56)} die op de \VUgras{tegelvloer} \textsuperscript{(57)} staat hebben
waarschijnlijk ook hun plaats in die kast.

%%K t06-c4-12-1 p
12. --- De keukenmeid heeft zeker geen gebrek aan fornuizen, want op de hoek van de tafel
staat ook nog een \VUgras{gasstel} \textsuperscript{(59)} waarop water in een kleine kastrol
warm wordt. De
\VUgras{stoom} \textsuperscript{(60)} die eruit ontsnapt wijst ons erop, dat het
waarschijnlijk kookt. Waarschijnlijk zal dat water voor de koffie gebruikt worden, want aan
het andere uiteinde van de tafel staan de
\VUgras{koffiekan} \textsuperscript{(64)} en de \VUgras{koffiemolen}
\textsuperscript{(63)}.

%%K t06-c4-13-1 p
13. --- Het stel is met de \VUgras{gasbrander} \textsuperscript{(62)} verbonden door een lang
\VUgras{rubberen slangetje} \textsuperscript{(61)}.

%%K t06-c4-14-1 p
14. --- De \VUgras{werkster} \textsuperscript{(66)} die de keukenmeid komt helpen, maakt de
groenten voor het diner klaar. Wanneer zij geschoond en gewassen zullen zijn, zal zij ze in
het \VUgras{bekken} \textsuperscript{(65)} leggen in afwachting van het uur om ze te koken.

%%K t06-tit-10 sub
\VUsaut{4.0mm}
{\centering\fontsize{10.2pt}{10.2pt}\selectfont\textls[40]{\scshape De badkamer.} \textit{(Badvertrek.)}\par}
\VUsaut{3.0mm}

%%K t06-c5-01-1 p
1. --- Ons rest nog maar één onbescheidenheid te begaan om de voornaamste kamers van de woning
te leren kennen: de inrichting van de badkamer te zien. Dat zal niet veel tijd vragen, want
zij is niet groot, en wij zullen snel weten dat de
\VUgras{badkuip} \textsuperscript{(1)} een goede \VUgras{geiser}
\textsuperscript{(2)} heeft, dat het \VUgras{zeepbakje} \textsuperscript{(3)} binnen het
bereik van de hand van de bader is en dat het
\VUgras{handdoekenrek} \textsuperscript{(5)} de \VUgras{handdoeken}
\textsuperscript{(4)} zeker gemakkelijk zal laten drogen.

%%K t06-c5-02-1 p
2. --- Die kamer heeft haar wanden bedekt met \VUgras{tegelwerk} \textsuperscript{(6)}, dat
het mogelijk maakte een \VUgras{douchetoestel} \textsuperscript{(7)} te plaatsen zonder
enigszins te vrezen dat het water, dat tijdens het gebruik ervan terugspat, de muren zou
beschadigen. Een zinken
\VUgras{teiltje} \textsuperscript{(8)} dat voor de voetbaden dient, maakt het
meubilair compleet.
\VUsaut{6.0mm}
\VUfilet{22mm}
